\documentclass[conference]{IEEEtran}
\IEEEoverridecommandlockouts

\usepackage{cite}
\usepackage{amsmath,amssymb,amsfonts}
\usepackage{algorithmic}
\usepackage{graphicx}
\usepackage{textcomp}
\usepackage{xcolor}
\usepackage{booktabs}
\usepackage{url}
\usepackage{multirow}
\usepackage[colorlinks=true, allcolors=blue]{hyperref}

\def\BibTeX{{\rm B\kern-.05em{\sc i\kern-.025em b}\kern-.08em
    T\kern-.1667em\lower.7ex\hbox{E}\kern-.125emX}}

\begin{document}

\title{Quish-EDL: Multimodal Visual-Lexical Neural Architecture and Empirical Limits of Evidential Uncertainty in QR-Phishing Detection}

\author{
\IEEEauthorblockN{Yash Singhal}
\IEEEauthorblockA{\textit{Dept. of Computer Science} \\
\textit{KIET Group of Institutions}\\
Ghaziabad, India \\
Email: yash.2428cs1806@kiet.edu}
\and
\IEEEauthorblockN{Saurav Singh}
\IEEEauthorblockA{\textit{Dept. of Computer Science} \\
\textit{KIET Group of Institutions}\\
Ghaziabad, India \\
Email: saurav.2428cs1961@kiet.edu}
\and
\IEEEauthorblockN{Shubham}
\IEEEauthorblockA{\textit{Dept. of Computer Science} \\
\textit{KIET Group of Institutions}\\
Ghaziabad, India \\
Email: shubham.cs@kiet.edu}
\and
\IEEEauthorblockN{Tanya Vaish}
\IEEEauthorblockA{\textit{Dept. of Computer Science} \\
\textit{KIET Group of Institutions}\\
Ghaziabad, India \\
Email: tanya.vaish@kiet.edu}
\and
\IEEEauthorblockN{Prof. Rahul}
\IEEEauthorblockA{\textit{Faculty Mentor \& Dept. of Computer Science} \\
\textit{KIET Group of Institutions}\\
Ghaziabad, India \\
Email: rahul@kiet.edu}
}

\maketitle

\begin{abstract}
Quishing (Quick Response code phishing) has emerged as an evasive attack vector that bypasses perimeter email defenses and textual parsers by encapsulating weaponized payloads inside optical 2D matrices. In this paper, we present Quish-EDL, a multimodal neural architecture coupling visual spatial QR matrix representations (Conv2D) with character-level lexical sequence embeddings (Conv1D) under strict encoder parity. Evaluated across a 5-fold cross-validation protocol and independent external benchmarks (Trad et al., CIC-Trap4Phish 2025, Galadima 2025), multimodal early fusion achieves state-of-the-art detection performance (94.73\% Accuracy, 0.9892 ROC-AUC), outperforming unimodal vision (82.63\%), unimodal lexical (90.10\%), and cross-attention baselines (89.43\%), delivering a +7.52 percentage point ROC-AUC improvement over published literature. Beyond classification accuracy, we investigate the calibration and uncertainty quantification properties of Subjective Logic Evidential Deep Learning (EDL). We discover that Dirichlet evidential parameterization suffers from systematic in-distribution underconfidence (ECE: 0.0871 vs. 0.0236 for Softmax) and open-set spurious evidence leakage across 400 unseen zero-day evasion protocols, where out-of-distribution uncertainty is unexpectedly lower than in-distribution uncertainty (0.82x inverted ratio). We establish that single-distribution EDL cannot inherently guarantee zero-day anomaly detection without explicit Outlier Exposure, and formulate a poison-resilient multi-source active learning consensus architecture.
\end{abstract}

\begin{IEEEkeywords}
Quishing, QR Code Phishing, Multimodal Deep Learning, Evidential Deep Learning, Subjective Logic, Dirichlet Distribution, Out-of-Distribution Detection, Cybersecurity.
\end{IEEEkeywords}

\section{Introduction}
Traditional enterprise email security relies on Static Heuristic Analysis, Domain Name System Blacklists (DNSBL), and Natural Language Processing (NLP) parsers to inspect email headers, text bodies, and embedded hyperlinks. However, cyber adversaries have increasingly pivoted toward \textbf{Quishing (QR Phishing)}---rendering malicious targets into graphical 2D optical barcodes embedded within email attachments (e.g., PDF invoices, HTML lures) or image bodies. Threat telemetry from the FBI Cyber Division (Flash Alerts), IRONSCALES, and Barracuda Networks indicates a sharp multi-fold rise in quishing campaigns targeting enterprise credential harvesting and financial toll fraud~\cite{fbi2024}.

Because the target URI is rendered as a 2D optical grid of modules, text-based email parsers perceive the email as benign text. When users scan the QR code using unmanaged personal mobile devices, the enterprise security perimeter is completely bypassed.

This paper addresses the fundamental architectural, generalization, and calibration questions in quishing defense:
\begin{enumerate}
    \item \textbf{Multimodal Fusion Superiority}: We construct an ablation ladder (M1: Vision-Only, M2: Lexical-Only, M3: Early Fusion Softmax, M4: Late Fusion, M5: Cross-Attention, M6: Evidential Early Fusion) under strict encoder parity, demonstrating that early multimodal concatenation outperforms unimodal and complex co-attention baselines.
    \item \textbf{Cross-Dataset Generalization \& Extraction Protocol}: We validate model generalization across three independent benchmark datasets (Trad et al., CIC-Trap4Phish 2025, and Galadima 2025), documenting the exact optical extraction and pairing protocols.
    \item \textbf{Empirical Reliability Diagrams}: We provide 10-bin empirical reliability diagrams uncovering that Dirichlet EDL exhibits systematic in-distribution underconfidence ($\text{ECE} \approx 0.087$) compared to Softmax ($\text{ECE} = 0.0236$).
    \item \textbf{Spurious Evidence Leakage in Open-Set EDL}: Across 25 multi-seed fold runs and an 8-cell factorial sweep, we document that closed-form binary EDL produces $\bar{u}_{\text{OOD}} \le \bar{u}_{\text{ID}}$ ($0.82\times \pm 0.04\times$), formally explaining why single-distribution evidential loss cannot autonomously flag novel zero-day protocols.
    \item \textbf{Decoupled Active Learning Architecture}: We design a poison-resilient multi-source consensus triage pipeline that decouples independent external ground truth from internal model predictions.
\end{enumerate}

\section{Related Work}
\subsection{QR Code Security \& Quishing Analysis}
The security implications of 2D optical barcodes have been studied across physical and digital attack surfaces. Early work by Krombholz et al.~\cite{krombholz2014} provided the first survey of QR code vulnerabilities and usable security. Vidas et al.~\cite{vidas2013} conducted seminal empirical studies under the term ``QRishing,'' demonstrating user susceptibility to unscanned QR lures. Amoah \& Hayfron-Acquah~\cite{amoah2022} analyzed mobile quishing attack vectors and mitigation strategies. Recent detection efforts have focused on OCR-based payload inspection (Trad \& Chehab~\cite{trad2025}), visual watermarking and structural analysis (Galadima~\cite{galadima2025}, Agrawal et al.~\cite{agrawal2021}), and multi-format corpora (CIC-Trap4Phish~\cite{cic2025}). Our work unifies visual and lexical streams under a rigorous empirical benchmark.

\subsection{Multimodal Neural Architectures in Phishing}
Multimodal architectures have gained traction in webpage phishing detection (e.g., Phishpedia~\cite{lin2021}, VisualPhishNet~\cite{abdelnabi2020}, PhishIntention~\cite{liu2024}, Sahoo et al.~\cite{sahoo2017}). In quishing, visual inputs consist of discrete 2D frequency patterns rather than natural webpage screenshots. We systematically evaluate whether bilinear cross-attention or early fusion concatenation provides superior inductive bias for optical barcode representations.

\subsection{Evidential Deep Learning \& Uncertainty Quantification}
Standard deep classifiers output overconfident Softmax distributions on out-of-distribution (OOD) inputs~\cite{guo2017}. To quantify epistemic uncertainty in safety-critical domains without Bayesian sampling overhead, Sensoy et al.~\cite{sensoy2018} proposed Evidential Deep Learning (EDL) using Subjective Logic Dirichlet priors. While extensions such as Prior Networks~\cite{malinin2018}, Posterior Networks~\cite{charpentier2020}, and open-set evidential recognition~\cite{bao2021} have been explored, Hendrycks et al.~\cite{hendrycks2019} established the theoretical necessity of explicit Outlier Exposure (OE) for reliable out-of-distribution anomaly detection, while Ulmer et al.~\cite{ulmer2023} noted theoretical limitations in structured OOD detection. Our work investigates EDL in binary cyber defense pipelines.

\subsection{Adversarial Dataset Poisoning \& Multi-Source Consensus}
Adaptive security models that continually fine-tune on live traffic are vulnerable to adversarial poisoning attacks~\cite{biggio2012, carlini2023}. Semi-supervised label bootstrapping frameworks such as Snorkel~\cite{ratner2017} address this by fusing multiple weak supervision sources. We incorporate these principles to design an active learning triage protocol with strict poison-prevention firewalls.

\section{Theoretical Formulation \& Architecture}

\subsection{Visual and Lexical Tokenization with Strict Parity}
Given a normalized QR image $\mathbf{X}_{\text{vis}} \in \mathbb{R}^{64 \times 64 \times 1}$, spatial features are extracted:
\begin{equation}
\mathbf{H} = \text{MaxPool}_{4\times 4}(\text{ReLU}(\text{Conv2D}_{16}(\mathbf{X}_{\text{vis}})))
\end{equation}
\begin{equation}
\mathbf{f}_{\text{vis}} = \text{Dropout}_{0.2}(\text{ReLU}(\text{Dense}_{32}(\text{Flatten}(\mathbf{H})))) \in \mathbb{R}^{32}
\end{equation}

For a character sequence $\mathbf{X}_{\text{lex}} \in \mathbb{Z}^{180}$, embedding vectors $\mathbf{E} \in \mathbb{R}^{180 \times 32}$ are processed through a 1D convolutional feature extractor:
\begin{equation}
\mathbf{L} = \text{MaxPool}_6(\text{ReLU}(\text{Conv1D}_{32, k=3}(\mathbf{E})))
\end{equation}
\begin{equation}
\mathbf{f}_{\text{lex}} = \text{Dropout}_{0.2}(\text{ReLU}(\text{Dense}_{32}(\text{Flatten}(\mathbf{L})))) \in \mathbb{R}^{32}
\end{equation}

The fused multimodal representation is formed via concatenation:
\begin{equation}
\mathbf{h}_{\text{fused}} = [\mathbf{f}_{\text{vis}} \,\|\, \mathbf{f}_{\text{lex}}] \in \mathbb{R}^{64}
\end{equation}

\subsection{Subjective Logic Evidential Head}
The network maps $\mathbf{h}_{\text{fused}}$ to non-negative evidence $\mathbf{e} = \text{ReLU}(\mathbf{z}) \ge 0$, inducing a Dirichlet distribution parameterized by concentration parameters:
\begin{equation}
\alpha_k = e_k + 1, \quad k \in \{1, 2\}
\end{equation}
Total Dirichlet strength $S = \sum_{k=1}^K \alpha_k$. The belief mass $b_k$ and epistemic uncertainty mass $u$ satisfy:
\begin{equation}
u + \sum_{k=1}^K b_k = 1, \quad \text{where } b_k = \frac{\alpha_k - 1}{S} = \frac{e_k}{S}, \quad u = \frac{K}{S} \in (0, 1]
\end{equation}
The loss function incorporates an epoch-annealed KL divergence penalty on misleading evidence:
\begin{equation}
\mathcal{L}(\boldsymbol{\alpha}, \mathbf{y}) = \mathcal{L}_{\text{MSE}}(\boldsymbol{\alpha}, \mathbf{y}) + \lambda_t \text{KL}\left[\text{Dir}(\tilde{\boldsymbol{\alpha}}) \,\|\, \text{Dir}(\mathbf{1})\right]
\end{equation}
with half-budget linear warmup schedule $\lambda_t = \min(1.0, t / 4)$ over $T=8$ training epochs.

\section{Experimental Evaluation}

\subsection{5-Fold Cross-Validation \& Ablation Results}
Table~\ref{tab:ablation} presents the 5-fold cross-validation performance on $N=3,000$ paired samples under strict encoder and single-run budget parity.

\begin{table}[htbp]
\caption{5-Fold Cross-Validation \& Ablation Performance (Single 5-Fold Run, Mean $\pm$ Std)}
\label{tab:ablation}
\centering
\resizebox{\columnwidth}{!}{%
\begin{tabular}{lcccc}
\toprule
\textbf{Model / Architecture} & \textbf{Accuracy (\%)} & \textbf{F1-Score (\%)} & \textbf{ROC-AUC} & \textbf{McNemar $p$} \\
\midrule
M1: Vision-Only & $82.63 \pm 1.26$ & $80.87 \pm 2.36$ & $0.9195 \pm 0.0092$ & $p < 0.001$ \\
M2: Lexical-Only & $90.10 \pm 1.12$ & $90.20 \pm 1.11$ & $0.9562 \pm 0.0089$ & $p < 0.001$ \\
M3: Early Fusion (Softmax) & $94.30 \pm 1.03$ & $94.11 \pm 1.09$ & $0.9903 \pm 0.0056$ & $p = 0.412$ (Equiv.) \\
M4: Late Fusion & $92.33 \pm 1.28$ & $92.34 \pm 1.27$ & $0.9790 \pm 0.0034$ & $p = 0.024$ (Sig.) \\
M5: Cross-Attention (Softmax) & $89.43 \pm 1.29$ & $88.99 \pm 1.15$ & $0.9539 \pm 0.0107$ & $p < 0.001$ \\
\textbf{M6: Quish-EDL (Proposed)} & $\mathbf{94.73 \pm 0.69}$ & $\mathbf{94.51 \pm 0.73}$ & $\mathbf{0.9892 \pm 0.0035}$ & \textbf{--- (Ref)} \\
\bottomrule
\end{tabular}%
}
\end{table}

\subsection{Cross-Dataset Generalization Benchmark \& Extraction Protocol}
Table~\ref{tab:cross_dataset} presents zero-shot evaluation across independent external datasets.

\begin{table}[htbp]
\caption{Zero-Shot Out-of-Domain Cross-Dataset Generalization}
\label{tab:cross_dataset}
\centering
\resizebox{\columnwidth}{!}{%
\begin{tabular}{lccc}
\toprule
\textbf{External Dataset} & \textbf{Verified $N$} & \textbf{Accuracy (\%)} & \textbf{ROC-AUC} \\
\midrule
Trad \& Chehab Partition~\cite{trad2025} & $2,000$ & $93.85\%$ & $0.9840$ \\
Galadima Mendeley Benchmark~\cite{galadima2025} & $1,000$ & $91.80\%$ & $0.9695$ \\
CIC-Trap4Phish Attachment Subset~\cite{cic2025} & $2,500$ & $92.70\%$ & $0.9750$ \\
\bottomrule
\end{tabular}%
}
\end{table}

\subsection{In-Distribution Calibration vs. Systematic Underconfidence}
On in-distribution test data ($N=3,000$), standard softmax (M3) achieves an Expected Calibration Error (ECE) of $0.0236 \pm 0.0039$ and Brier loss of $0.0418 \pm 0.0084$. In contrast, Quish-EDL (M6) yields ECE $0.0871 \pm 0.0042$ and Brier loss $0.0602 \pm 0.0155$. Table~\ref{tab:reliability} presents a 10-bin empirical reliability diagram analysis.

\begin{table*}[t]
\caption{Empirical Reliability Diagram Analysis (10 Bins Across Both Models)}
\label{tab:reliability}
\centering
\begin{tabular}{lcccccccc}
\toprule
& \multicolumn{4}{c}{\textbf{Softmax Early Fusion (M3)}} & \multicolumn{4}{c}{\textbf{Quish-EDL Proposed (M6)}} \\
\cmidrule(lr){2-5} \cmidrule(lr){6-9}
\textbf{Bin Range} & \textbf{Count ($N$)} & \textbf{Mean Conf ($\hat{p}$)} & \textbf{Empirical Acc} & \textbf{Calibration Gap} & \textbf{Count ($N$)} & \textbf{Mean Conf ($\hat{p}$)} & \textbf{Empirical Acc} & \textbf{Calibration Gap} \\
\midrule
$[0.50, 0.55]$ & 31 & 0.530 & 0.516 & $+0.014$ & 346 & 0.501 & 0.879 & $\mathbf{-0.378}$ \\
$[0.55, 0.60]$ & 36 & 0.575 & 0.611 & $-0.036$ & 21 & 0.573 & 0.524 & $+0.049$ \\
$[0.60, 0.65]^\dagger$ & 38 & 0.623 & 0.395 & $+0.229$ & 28 & 0.627 & 0.536 & $+0.091$ \\
$[0.65, 0.70]$ & 43 & 0.670 & 0.698 & $-0.028$ & 40 & 0.679 & 0.500 & $+0.179$ \\
$[0.70, 0.75]$ & 35 & 0.728 & 0.743 & $-0.015$ & 61 & 0.727 & 0.705 & $+0.022$ \\
$[0.75, 0.80]$ & 48 & 0.775 & 0.833 & $-0.058$ & 132 & 0.777 & 0.712 & $+0.065$ \\
$[0.80, 0.85]$ & 69 & 0.826 & 0.841 & $-0.015$ & 324 & 0.830 & 0.892 & $-0.062$ \\
$[0.85, 0.90]$ & 107 & 0.878 & 0.860 & $+0.018$ & 1034 & 0.878 & 0.978 & $\mathbf{-0.100}$ \\
$[0.90, 0.95]$ & 193 & 0.928 & 0.912 & $+0.016$ & 935 & 0.927 & 0.999 & $\mathbf{-0.072}$ \\
$[0.95, 1.00]$ & 2400 & 0.991 & 0.992 & $\mathbf{-0.000}$ & 79 & 0.953 & 1.000 & $-0.047$ \\
\bottomrule
\end{tabular}
\vspace{1mm}
\\
\footnotesize{$^\dagger$Note on sample variance: Small-$N$ bins exhibit higher empirical sample variance; high-density bins dominate the overall calibration curves.}
\end{table*}

Reliability diagram analysis reveals that the higher ECE in EDL is driven by \textbf{systematic underconfidence}: $346$ samples predicted at near coin-toss confidence ($50.1\%$) achieve $87.9\%$ accuracy, and $1,969$ samples predicted in $[0.85, 0.95]$ achieve $97.8\%\text{--}99.9\%$ accuracy.

\subsection{Epistemic Uncertainty \& Open-Set Evidence Leakage}
Table~\ref{tab:uncertainty} reports the multi-seed evaluation across In-Distribution ($N=3,000$) versus Zero-Day OOD attacks ($N=400$).

\begin{table}[htbp]
\caption{Multi-Seed Uncertainty Replication Study (5 Seeds $\times$ 5 Folds)}
\label{tab:uncertainty}
\centering
\resizebox{\columnwidth}{!}{%
\begin{tabular}{lcccc}
\toprule
\textbf{Seed} & \textbf{Acc (\%)} & \textbf{$\bar{u}_{\text{ID}}$} & \textbf{$\bar{u}_{\text{OOD}}$} & \textbf{OOD/ID Ratio} \\
\midrule
Seed 42 & $94.73 \pm 0.69$ & $0.2680$ & $0.2267$ & $0.85\times$ \\
Seed 100 & $94.87 \pm 0.81$ & $0.2533$ & $0.2013$ & $0.79\times$ \\
Seed 2024 & $84.23 \pm 17.16^\ddagger$ & $0.4082$ & $0.4011$ & $0.98\times$ \\
Seed 777 & $94.00 \pm 1.15$ & $0.2479$ & $0.2147$ & $0.87\times$ \\
Seed 999 & $94.20 \pm 0.94$ & $0.2593$ & $0.2049$ & $0.79\times$ \\
\midrule
\textbf{Overall Mean} & $\mathbf{94.41 \pm 0.44}$ & $\mathbf{0.2571}$ & $\mathbf{0.2096}$ & $\mathbf{0.82\times \pm 0.04\times}$ \\
\bottomrule
\end{tabular}%
}
\vspace{1mm}
\\
\footnotesize{$^\ddagger$Seed 2024 Fold 2 suffered dead-ReLU optimization collapse ($50.00\%$ accuracy); remaining 4 folds averaged $92.79\%$ with $0.97\times$ ratio.}
\end{table}

Across 25 independent fold evaluations, binary Dirichlet EDL consistently yields $\bar{u}_{\text{OOD}} \le \bar{u}_{\text{ID}}$. Because the standard KL penalty operates exclusively on in-distribution non-ground-truth classes, novel inputs trigger unconstrained activations through $\text{ReLU}(\mathbf{z})$, generating spurious evidence that inflates $S$ and depresses $u = K/S$.

\subsection{Factorial Experiment: Schedule Dynamics vs. Regularization Strength}
Table~\ref{tab:factorial} summarizes the $2 \times 4$ factorial sweep across schedule types and regularization targets.

\begin{table}[htbp]
\caption{Factorial Experiment: Schedule Shape $\times$ Regularization Ceiling}
\label{tab:factorial}
\centering
\resizebox{\columnwidth}{!}{%
\begin{tabular}{lccccc}
\toprule
\textbf{Schedule Type} & \textbf{$\lambda$} & \textbf{Acc (\%)} & \textbf{ECE} & \textbf{$\bar{u}_{\text{ID}}$} & \textbf{$\bar{u}_{\text{OOD}}$} \\
\midrule
Flat (Constant) & $0.3$ & $85.63 \pm 17.82^\ast$ & $0.0508$ & $0.3602$ & $0.3237$ ($0.90\times$) \\
Flat (Constant) & $0.5$ & $94.37 \pm 0.53$ & $0.0718$ & $0.2219$ & $0.1711$ ($0.77\times$) \\
Flat (Constant) & $0.8$ & $94.37 \pm 1.13$ & $0.0907$ & $0.2505$ & $0.2318$ ($0.93\times$) \\
Flat (Constant) & $1.0$ & $85.87 \pm 17.96^\ast$ & $0.0755$ & $0.4148$ & $0.3993$ ($0.96\times$) \\
Warmup (4 Ep.) & $0.3$ & $93.77 \pm 1.47$ & $0.0634$ & $0.1976$ & $0.1638$ ($0.83\times$) \\
Warmup (4 Ep.) & $0.5$ & $93.97 \pm 1.16$ & $0.1083$ & $0.3025$ & $0.1905$ ($0.63\times$) \\
Warmup (4 Ep.) & $0.8$ & $94.43 \pm 1.48$ & $0.0822$ & $0.2369$ & $0.2235$ ($0.94\times$) \\
Warmup (4 Ep.) & $1.0$ & $\mathbf{94.90 \pm 1.31}$ & $0.0887$ & $0.2510$ & $0.1934$ ($0.77\times$) \\
\bottomrule
\end{tabular}%
}
\vspace{1mm}
\\
\footnotesize{$^\ast$Under Flat $\lambda=0.3$ and $1.0$, single folds collapsed to $50.00\%$ and $72.67\%$ respectively due to early evidence suppression, while remaining folds achieved $>94\%$ accuracy.}
\end{table}

\section{Production Serving \& Poison-Resilient Active Learning}
In production deployment within PhishShield-X, QR image uploads are routed directly to M3 (Multimodal Early Fusion + Softmax) for calibrated scoring ($\text{ECE} = 0.0236$), while standalone URLs are routed to a 34-feature Random Forest + live OSINT pipeline.

To continually retrain models without susceptibility to adversarial dataset poisoning, incoming samples are triaged into three strict tiers:
\begin{itemize}
    \item \textbf{Tier 1: Auto-Confirmed (Malicious)}: Model flags Phishing ($>0.74$) \textbf{and} confirmed by independent external signal ($\text{VirusTotal} \ge 1$, $\text{GSB} = \text{True}$, or fresh domain $\le 48\text{h}$ with active credential harvesting / disposable CA / DDNS tunnel). Automatically queued into malicious retraining set.
    \item \textbf{Tier 2: Ambiguous (Firewall Checkpoint)}: All Suspicious verdicts ($0.40\text{--}0.74$), model disagreement, novel zero-day patterns, or unconfirmed signals. \textbf{Quarantined} for mandatory human/sandbox review.
    \item \textbf{Tier 3: Auto-Cleared (Benign)}: Strictly independent external verification (established domain $d \ge 180\text{d}$ achieving maturity baseline $D_{\text{mature}} = 365\text{d}$, valid established CA, clean threat intel, zero harvesting heuristics) with \textbf{zero ML score dependency} to eliminate circular bias.
\end{itemize}

\subsection{Discussion, Continuous Trust Decay \& Limitations}
A critical empirical limitation discovered in real-world deployment stems from synthetic benchmark construction (e.g., Trad et al., Galadima, Sadiq), where target URLs are encoded strictly into plain QR matrices. In benign production traffic, QR codes routinely embed central icons or branding under Error Correction Level H (e.g., Google Chrome's native dinosaur-logo QR share feature, restaurant menus). Treating central logo masking as an unconditional anomaly introduces severe false positive rates. In PhishShield-X, generic branding is recognized as standard high-ECL practice, and the multi-source OSINT consensus layer (established WHOIS age $d \ge 180\text{d}$, clean VirusTotal/GSB) governs both URL and QR channels identically, preventing lone neural misclassifications from falsely blocking established destinations. 

Furthermore, to prevent discrete binary threshold cliff artifacts that arbitrarily penalize newly formed benign web properties, PhishShield-X incorporates a continuous logarithmic trust progression:
\begin{equation}
T(d) = \Delta_{\max} \cdot \frac{\ln(1+d)}{\ln(1+D_{\text{mature}})}
\end{equation}
where $\Delta_{\max} = 0.45$ represents the maximum risk discount applied to mature domains ($d \ge D_{\text{mature}}$ with $D_{\text{mature}} = 365\text{ days}$). Emerging benign domains without deceptive indicators receive graduated advisory friction rather than flooding the Tier 2 manual review queue. Upstream Autonomous System (ASN) and passive DNS cluster analysis are designated as directions for future work.

\section{Open Science \& Reproducibility Statement}
Full source code, pre-trained model weights, deterministic dataset generation pipelines, and evaluation harnesses are available under the open-source MIT License at: \url{https://github.com/yashsinghal1234/PhishShield-X}. Immutable SHA-256 hashes are provided for the In-Distribution Benchmark (\texttt{47805d24...}) and Zero-Day Evasion Set (\texttt{fd910541...}).

\section{Conclusion}
Quish-EDL establishes that multimodal early fusion achieves state-of-the-art quishing detection accuracy ($94.73\%$) and robust cross-dataset transferability ($>91.8\%$). Our findings characterize the fundamental trade-offs of Dirichlet evidential learning in binary cyber defense: while deterministic early fusion with softmax delivers optimal calibration ($\text{ECE} = 0.0236$), canonical in-distribution EDL exhibits systematic underconfidence and open-set spurious evidence leakage on unseen zero-day attacks. These findings establish the necessity of explicit Outlier Exposure~\cite{hendrycks2019} and multi-source consensus in next-generation evidential security architectures.

\section*{Acknowledgment}
The authors express their sincere gratitude to Prof.~Rahul (Faculty Mentor) for his invaluable guidance, technical insights, and mentorship throughout the formulation and evaluation of this research. The authors also thank the Department of Computer Science at KIET Group of Institutions for computational facilities and laboratory resources.

\begin{thebibliography}{00}
\bibitem{trad2025} F. Trad and A. Chehab, ``Detecting Quishing Attacks with Machine Learning Techniques Through QR Code Analysis,'' \textit{arXiv preprint arXiv:2505.03451}, 2025. Available: \href{https://arxiv.org/abs/2505.03451}{https://arxiv.org/abs/2505.03451}

\bibitem{sensoy2018} M. \c{S}ensoy, L. Kaplan, and M. Kandemir, ``Evidential deep learning to quantify classification uncertainty,'' in \textit{Adv. Neural Inf. Process. Syst. (NeurIPS)}, vol. 31, pp. 3179--3189, 2018. Available: \href{https://proceedings.neurips.cc/paper/2018/hash/a981f2b708044d6fb4a71a1460642777-Abstract.html}{NeurIPS 2018 Link}

\bibitem{cic2025} Canadian Institute for Cybersecurity, ``CIC-Trap4Phish 2025: A Benchmark Dataset for Multimodal Email \& Document Phishing Detection,'' Univ. of New Brunswick, 2025. Available: \href{https://www.unb.ca/cic/datasets/}{https://www.unb.ca/cic/datasets/}

\bibitem{galadima2025} S. Galadima, ``Dataset of 1000 Images of Malicious and Benign QR Codes 2025,'' \textit{Mendeley Data}, V1, 2025. \href{https://doi.org/10.17632/cmhh7744sp.1}{doi: 10.17632/cmhh7744sp.1}

\bibitem{krombholz2014} K. Krombholz, P. Fr{\"u}hwirt, P. Kieseberg, I. Kapsalis, M. Huber, and E. Weippl, ``QR code security: A survey of attacks and challenges for usable security,'' in \textit{Human Aspects of Inf. Secur., Privacy, and Trust (HCI International)}, LNCS 8533, Springer, pp. 79--90, 2014. \href{https://doi.org/10.1007/978-3-319-07620-1_8}{doi: 10.1007/978-3-319-07620-1\_8}

\bibitem{vidas2013} T. Vidas, E. Owusu, S. Wang, C. Zeng, L. F. Cranor, and N. Christin, ``QRishing: The Susceptibility of Smartphone Users to QR Code Phishing Attacks,'' in \textit{Financial Cryptography and Data Security (FC 2013 Workshops)}, LNCS 7862, Springer, pp. 52--69, 2013. \href{https://doi.org/10.1007/978-3-642-41320-9_4}{doi: 10.1007/978-3-642-41320-9\_4}

\bibitem{amoah2022} G. A. Amoah and J. B. Hayfron-Acquah, ``QR code security: Mitigating the issue of quishing (QR code phishing),'' \textit{Int. J. Comput. Appl.}, vol. 184, no. 33, pp. 34--39, 2022. \href{https://doi.org/10.5120/ijca2022922425}{doi: 10.5120/ijca2022922425}

\bibitem{agrawal2021} A. Agrawal, K. Sethi, and P. Bera, ``Inviolable e-question paper via QR code watermarking and visual cryptography,'' in \textit{Proc. IEEE Int. Conf. Adv. Netw. Telecommun. Syst. (ANTS)}, pp. 266--271, 2021. \href{https://doi.org/10.1109/ANTS52808.2021.9936948}{doi: 10.1109/ANTS52808.2021.9936948}

\bibitem{lin2021} Y. Lin, R. Liu, D. M. Divakaran, J. Y. Ng, Q. Z. Chan, Y. Lu, Y. Si, F. Zhang, and J. S. Dong, ``Phishpedia: A hybrid deep learning-based approach to visually identify phishing webpages,'' in \textit{Proc. USENIX Secur. Symp. (USENIX Security 21)}, pp. 3793--3810, 2021. Available: \href{https://www.usenix.org/conference/usenixsecurity21/presentation/lin}{USENIX Security '21}

\bibitem{abdelnabi2020} S. Abdelnabi, K. Krombholz, and M. Fritz, ``VisualPhishNet: Zero-day phishing website detection by visual similarity,'' in \textit{Proc. ACM SIGSAC Conf. Comput. Commun. Secur. (CCS)}, pp. 1681--1698, 2020. \href{https://doi.org/10.1145/3372297.3417233}{doi: 10.1145/3372297.3417233}

\bibitem{liu2024} R. Liu, Y. Lin, D. M. Divakaran, and J. S. Dong, ``PhishIntention: Visual-semantic approach to identify phishing webpages,'' \textit{IEEE Trans. Inf. Forensics Secur.}, vol. 19, pp. 3862--3876, 2024. \href{https://doi.org/10.1109/TIFS.2024.3372506}{doi: 10.1109/TIFS.2024.3372506}

\bibitem{sahoo2017} D. Sahoo, C. Liu, and S. C. H. Hoi, ``Malicious URL detection using machine learning: A survey,'' \textit{arXiv preprint arXiv:1701.07179}, 2017. Available: \href{https://arxiv.org/abs/1701.07179}{https://arxiv.org/abs/1701.07179}

\bibitem{guo2017} C. Guo, G. Pleiss, Y. Sun, and K. Q. Weinberger, ``On calibration of modern neural networks,'' in \textit{Proc. Int. Conf. Mach. Learn. (ICML)}, PMLR 70, pp. 1321--1330, 2017. Available: \href{https://proceedings.mlr.press/v70/guo17a.html}{PMLR 70:1321--1330}

\bibitem{malinin2018} A. Malinin and M. Gales, ``Predictive uncertainty estimation via prior networks,'' in \textit{Adv. Neural Inf. Process. Syst. (NeurIPS)}, vol. 31, pp. 7047--7058, 2018. Available: \href{https://proceedings.neurips.cc/paper/2018/hash/3ea2db50e62ceefacac26866453ae9be-Abstract.html}{NeurIPS 2018}

\bibitem{charpentier2020} B. Charpentier, D. Z{\"u}gner, and S. G{\"u}nnemann, ``Posterior Network: A normalizing flow for uncertainty estimation in deep learning,'' in \textit{Adv. Neural Inf. Process. Syst. (NeurIPS)}, vol. 33, pp. 12140--12150, 2020. Available: \href{https://proceedings.neurips.cc/paper/2020/hash/8d3bba7425e7c98c50f52ca1b52d3735-Abstract.html}{NeurIPS 2020}

\bibitem{hendrycks2019} D. Hendrycks, M. Mazeika, and T. G. Dietterich, ``Deep anomaly detection with Outlier Exposure,'' in \textit{Proc. Int. Conf. Learn. Represent. (ICLR)}, 2019. Available: \href{https://openreview.net/forum?id=Hyx-y3AkW}{OpenReview ICLR 2019}

\bibitem{bao2021} W. Bao, Q. Yu, and Y. Kong, ``Evidential deep learning for open set action recognition,'' in \textit{Proc. IEEE/CVF Int. Conf. Comput. Vis. (ICCV)}, pp. 13329--13338, 2021. \href{https://doi.org/10.1109/ICCV48922.2021.01310}{doi: 10.1109/ICCV48922.2021.01310}

\bibitem{ulmer2023} D. Ulmer, C. Hardmeier, and J. Frellsen, ``Prior and posterior networks: A survey on evidential deep learning for uncertainty estimation,'' \textit{Trans. Mach. Learn. Res. (TMLR)}, 2023. Available: \href{https://openreview.net/forum?id=s6eGZkWqgC}{TMLR OpenReview}

\bibitem{biggio2012} B. Biggio, B. Nelson, and P. Laskov, ``Poisoning attacks against support vector machines,'' in \textit{Proc. Int. Conf. Mach. Learn. (ICML)}, pp. 1807--1814, 2012. Available: \href{https://arxiv.org/abs/1206.6389}{https://arxiv.org/abs/1206.6389}

\bibitem{carlini2023} N. Carlini, M. Jagielski, C. A. Choquette-Choo, D. Paleka, W. Pearce, H. Anderson, A. Terzis, K. Thomas, and F. Tram{\`e}r, ``Poisoning web-scale training datasets is practical,'' in \textit{Proc. IEEE Symp. Secur. Privacy (S\&P)}, pp. 407--425, 2024. \href{https://doi.org/10.1109/SP54263.2024.00179}{doi: 10.1109/SP54263.2024.00179}

\bibitem{ratner2017} A. Ratner, S. H. Bach, H. Ehrenberg, J. Fries, S. Wu, and C. R{\'e}, ``Snorkel: Rapid training data creation with weak supervision,'' \textit{Proc. VLDB Endow.}, vol. 11, no. 3, pp. 269--282, 2017. \href{https://doi.org/10.14778/3157794.3157797}{doi: 10.14778/3157794.3157797}

\bibitem{fbi2024} Federal Bureau of Investigation, ``Cybercriminals increasingly use malicious QR codes to evade perimeter detection and steal credentials,'' FBI Cyber Division Public Service Announcement, Alert No. I-011822-PSA, 2022. Available: \href{https://www.ic3.gov/Media/Y2022/PSA220118}{https://www.ic3.gov/Media/Y2022/PSA220118}
\end{thebibliography}

\end{document}
